\scanpage{36}

\unsure{L'opération|L'application|Définition} $\sigma_e$ s'appelle \uline{symétrie} \uline{p.r.} à la \uline{droite} $D$, comme $\sigma$ est un autom. de $A$ pour la structure quad. définie par la norme, $\sigma_e$ est un autom. de la structure quadratique de $E_e$ pour $Q$, a fortiori un autom. de la structure de similitude.

\uline{Prop} Supposons qu'en tout point de $X_0=V(2.1_{\underline O_X})$, i.e. tout pt de car. rés. $2$, la structure de similitude donnée sur $E$ soit non dég. (i.e $A$ soit étale sur $\underline O$ en ce point). Alors pour tte droite \deleted{tot}alement non isotrope $D\subset E$, $\sigma_D$ est l'\uline{unique} autom. $\Theta$ de la structure de similitude $\Sigma$ (ou encore, de la structure quadratique, quand $\Sigma$ est donnée par une forme quadratique $Q$) \deleted{telle} qui induise $\mathrm{id}_D$ sur $D$, $-\mathrm{id}$ sur $E/D$,$^*$ et qui \inserted{\unsure{ne|en}} soit \inserted{l'identité au voisin.} \deleted{\unsure{en} \unsure{les} points $s$ $\unsure{\in}S$} \deleted{\unsure{d'étale} de chaque pt $s$ de $X$}\deleted{--- est \unsure{induit},} \inserted{\unsure{en les|qu'en} points $s$ de $X-X_0$)}

% slanted
\margin{Ceci équivaut à $\det\Theta=-1$ (compte tenu que $\Theta|D=\mathrm{id}_D$)}

\uline{NB} Cette dernière condition \unsure{bien|n'est} \unsure{sûre.|nécessaire} est \unsure{auto}

\uline{Cor.} Soit $A$ algèbre quadratique sur $X$, supposons qu'en tt $x\in X_0$, $A$ soit étale en $x$. Alors l'involution canonique est le seul automorphisme \inserted{d'Alg} de $A$ qui induise $-\mathrm{id}.$ sur $A/\underline O.1_A$.

% archive sheet number: 35
